% !TEX root = ../my-thesis.tex

\lstset{breaklines=true, breakatwhitespace=false}

\chapter{Methodik}
\label{sec:methodik}

Dieses Kapitel beschreibt das Vorgehen der Untersuchung. Es legt das
Evaluationsdesign fest, begründet die Auswahl der fünf Workflow-Management-Systeme, führt die qualitativen Vergleichsaspekte und die quantitativen
Messgrößen ein und beschreibt die Benchmark-Workflows, den
Versuchsplan sowie das Mess- und Auswertungsverfahren.

\section{Evaluationsdesign}
\label{sec:methodik:design}

Die Untersuchung vergleicht fünf Workflow-Management-Systeme entlang
zweier Workflow-Muster und verbindet dabei zwei Blickrichtungen.
Qualitativ wird zunächst betrachtet, wie sich die Muster in den
Systemen beschreiben lassen und wie sich die Umsetzungen entlang
fester Vergleichsaspekte unterscheiden. Anschließend wird
quantitativ gemessen, wie lange die Systeme für die Ausführung der
Muster benötigen, welchen Zeitanteil die Systeme der reinen
Rechenlogik hinzufügen, wie sie die Übergaben zwischen den Schritten
abwickeln und wie die parallelen Instanzen gestartet werden.

Die Untersuchung ist zweistufig aufgebaut. Auf der lokalen Ebene
werden beide Muster in allen fünf Systemen auf einem einzelnen
Rechner umgesetzt und vermessen. Diese Ebene bildet die
Vergleichsbasis, da dort alle fünf Systeme vertreten sind und
unter denselben Bedingungen ausgeführt werden. Auf der zweiten Ebene wird untersucht, wie sich die Workflow-Management-Systeme auf einem Slurm-basierten HPC-System bereitstellen und
betreiben lassen, ergänzt um Messungen auf einer dediziert
zugewiesenen Ressource dieses Systems. Die lokale Ebene behandelt
Kapitel~\ref{sec:lokal}, die HPC-Ebene Kapitel~\ref{sec:mogon}.

In allen Systemen kommen dieselben Python-Skripte als Rechenschritte
zum Einsatz. Unterschiede in den gemessenen Laufzeiten gehen damit
nicht auf unterschiedliche fachliche Rechenoperationen zurück.

\section{Auswahl der Systeme}
\label{sec:methodik:auswahl}
Untersucht werden Nextflow, StreamFlow, Merlin, Pegasus und AiiDA. Alle fünf sind für wissenschaftliche Workflows entwickelt, unterscheiden sich aber in ihrem Entstehungskontext. Nextflow stammt aus dem Umfeld der Bioinformatik und zielt auf die reproduzierbare Ausführung von Analysen großer biologischer Datensätze über unterschiedliche Rechenumgebungen hinweg~\cite[S.~316]{ditommaso2017nextflow}. StreamFlow zielt darauf, denselben Workflow in verschiedenen Ausführungsumgebungen auszuführen~\cite[S.~1723]{colonnelli2021streamflow}. Merlin ist für große Ensembles wissenschaftlicher Simulationen ausgelegt, deren Ergebnisse für maschinelles Lernen aufbereitet werden~\cite[S.~255]{peterson2022merlin}. Pegasus richtet sich auf großskalige, mehrstufige Simulations- und Analyse-Abläufe, deren Rechen- und Datenumfang ein System erfordert, das Datentransfer und Aufgabenausführung auf verteilten Rechenressourcen zuverlässig koordiniert und automatisiert~\cite[S.~17]{deelman2015pegasus}. AiiDA verbindet die Ausführung mit einer durchgängigen Aufzeichnung der Datenherkunft in einer Datenbank~\cite[S.~1]{huber2020aiida}. Eine Kurzvorstellung der Systeme gibt Kapitel~\ref{sec:basics}.

Die Auswahl zielte darauf, unterschiedliche Ansätze
wissenschaftlicher Workflow-Management-Systeme abzudecken, Grundlage
waren die Veröffentlichungen und Dokumentationen der Systeme.
Berücksichtigt wurden drei Merkmale. Das erste ist der Anwendungskontext, in dem die Systeme entstanden sind. Dieser wurde oben je System benannt und umfasst unter anderem Bioinformatik, große Simulationsensembles und einen besonderen Schwerpunkt auf Provenance. Auch Suter
et al. unterscheiden Systeme danach, ob sie einer Fachgemeinschaft
entstammen oder anwendungsneutral ausgelegt
sind~\cite[S.~4]{suter2026terminology}. Das zweite
Merkmal ist die Art der Workflow-Beschreibung, die Systeme
unterscheiden sich darin, in welcher Form ein Workflow
niedergeschrieben wird und wie darin die Abhängigkeiten zwischen den
Schritten und die parallele Verarbeitung ausgedrückt werden. Das
dritte Merkmal ist die Architektur der Ausführung, einige Systeme
führen Workflows in einer eigenen Laufzeit aus, andere übergeben die
Ausführung an eigenständige Komponenten. Wie sich diese Merkmale in den
einzelnen Systemen zeigen, zeigt die Modellierung in
Kapitel~\ref{sec:lokal}.

Hinzu kamen praktische Erwägungen. Nextflow steht auf dem
betrachteten HPC-System bereits als Modul zur Verfügung, StreamFlow
wurde von der betreuenden Arbeitsgruppe vorgeschlagen.
Für alle Systeme wurde vorausgesetzt, dass sie frei verfügbar sind, aktiv
weiterentwickelt werden~\cite{wcisystems}, sich auf einem einzelnen
Rechner installieren lassen und eine öffentliche Dokumentation
besitzen, da beide Muster zunächst lokal umgesetzt und vermessen
werden.

Das Feld wissenschaftlicher Workflow-Management-Systeme umfasst
mehrere hundert Vertreter~\cite{suter2026terminology,wcisystems}.
Eine vollständige Erhebung ist im Rahmen dieser Arbeit nicht
möglich, da jedes System einzeln bereitgestellt, in beiden Mustern
umgesetzt und vermessen wird. Die Auswahl ist deshalb eine
begründete Stichprobe unterschiedlicher Ansätze.

\section{Qualitative Vergleichsaspekte}
\label{sec:methodik:qualitativ}
Der qualitative Vergleich betrachtet die Umsetzungen entlang von
sieben Aspekten: dem Ausdruck des Patterns, dem Modellierungs- und
Konfigurationsaufwand, der Sichtbarkeit von Datenfluss und Abhängigkeiten, der
Ausführung und Steuerung, der Trennung von Workflow-Logik und
Ausführungsumgebung, der Nachvollziehbarkeit der Ausführung 
und der Verständlichkeit aus Einsteigerperspektive. Zusammen decken sie den
Weg eines Workflows von der Beschreibung über die Ausführung bis
zur Nachvollziehbarkeit ab. Die folgenden Absätze
begründen die Wahl der einzelnen Aspekte.

\emph{Ausdruck des Patterns:} Dieser Aspekt betrachtet, mit welchen Mitteln eines Systems die Struktur der beiden Workflow-Muster beschrieben wird, insbesondere die Abfolge der Schritte und die parallele Verarbeitung der Teilstücke. Er bildet die Grundlage des Vergleichs, da die Workflow-Muster in allen fünf Systemen dieselben sind, ihre Umsetzung jedoch unterschiedlich erfolgt. Untersucht wird der Aspekt anhand der Workflow-Beschreibungen beider Muster in allen fünf Systemen.

\emph{Modellierungs- und Konfigurationsaufwand:} Dieser Aspekt betrachtet, welcher Beschreibungs- und Konfigurationsaufwand nötig ist, um ein Workflow-Muster in einem System umzusetzen und anzupassen. Da sich der Aufwand nicht unmittelbar messen lässt, wird er über zwei nachvollziehbare Größen erfasst: die Zahl der manuell erstellten Artefakte zur Beschreibung und Konfiguration des
Workflows sowie die Frage, welche Stellen bei einer Änderung der Rechenlast oder der Teilstückzahl angepasst werden müssen. Beide Größen sind unmittelbar an den Workflow-Beschreibungen und Konfigurationsdateien der Umsetzungen ablesbar.

\emph{Sichtbarkeit von Datenfluss und Abhängigkeiten:} Dieser Aspekt betrachtet, ob sich aus der Workflow-Beschreibung ablesen lässt, welche Dateien zwischen den Schritten weitergegeben werden und wodurch die Abhängigkeiten zwischen ihnen entstehen. Je unmittelbarer diese Informationen aus der Beschreibung hervorgehen, desto leichter lässt sich der Ablauf eines Workflows nachvollziehen. Der
Aspekt wird anhand der Workflow-Beschreibungen beider Muster in
allen fünf Systemen untersucht.

\emph{Ausführung und Steuerung:} Dieser Aspekt betrachtet, welche
Komponente einen beschriebenen Workflow ausführt und wodurch der
Beginn eines einzelnen Schritts ausgelöst wird. Damit bildet er die
Verbindung zwischen der Workflow-Beschreibung und dem beobachtbaren
Ablauf eines Workflows, da sich daraus ergibt, wann und auf welchem
Weg einzelne Schritte ausgeführt werden. Der Aspekt wird anhand der
Umsetzung und Ausführung der beiden Benchmark-Workflows betrachtet.

\emph{Trennung von Workflow-Logik und Ausführungsumgebung:} Dieser Aspekt betrachtet, ob die fachliche Beschreibung eines Workflows unabhängig von der Umgebung formuliert ist, in der er ausgeführt wird, und über welchen Mechanismus beides verbunden wird. Eine Trennung von Workflow-Logik und Ausführungsumgebung erleichtert es, denselben Workflow in unterschiedlichen Ausführungsumgebungen einzusetzen und ist für diese Arbeit besonders relevant, da beide Workflow-Muster zunächst lokal und anschließend auf einem HPC-System betrachtet werden. Mehrere der untersuchten Systeme nennen diese Trennung selbst als Entwurfsziel~\cite[S.~1723]{colonnelli2021streamflow}\cite[S.~18--20]{deelman2015pegasus}. Untersucht wird der Aspekt anhand der Mechanismen, mit denen die Systeme Workflow-Beschreibung und Ausführungsumgebung voneinander trennen, sowie daran, welche Teile der Umsetzung bei einem Wechsel der Ausführungsumgebung angepasst werden müssten.

\emph{Nachvollziehbarkeit der Ausführung:} Dieser Aspekt betrachtet, welche Möglichkeiten ein System bietet, den Fortschritt eines Laufs zu verfolgen, und welche Artefakte nach dem Lauf zur Verfügung stehen. Die Aufzeichnung des Ausführungsverlaufs erleichtert es, Ergebnisse nachzuvollziehen und Fehlerquellen im Ablauf einzugrenzen und ist insbesondere in komplexen Workflows und HPC-Umgebungen von Bedeutung~\cite{davidson2008provenance,sacco2024provenance}. Untersucht wird der Aspekt anhand der Ausgaben während der Ausführung, der verbleibenden Verzeichnisse und Protokolle sowie weiterer Artefakte und Darstellungen, die sich aus einem Lauf erzeugen lassen.

\emph{Verständlichkeit aus Einsteigerperspektive:} Dieser Aspekt
betrachtet, welche systemeigenen Konzepte verstanden sein müssen,
bevor sich ein Workflow in einem System umsetzen und ausführen
lässt. Die Einstiegshürde ist für die Wahl eines Systems von
praktischer Bedeutung.  Untersucht wird der Aspekt anhand der Konzepte, die für die erstmalige Umsetzung der beiden Workflow-Muster in den einzelnen Systemen verstanden werden mussten. Die Einordnung beruht auf dieser erstmaligen Umsetzung und ersetzt keine Nutzerstudie.


\section{Quantitative Messgrößen}
\label{sec:methodik:messgroessen}
Der quantitative Vergleich stützt sich auf vier Messgrößen: die Gesamtlaufzeit (Makespan), den Overhead gegenüber einem Referenzlauf sowie die beiden Koordinationsgrößen Übergangslatenz und Startspreizung. Die folgenden Absätze definieren die vier
Größen.

\emph{Gesamtlaufzeit:} Als Gesamtlaufzeit wird die Zeitspanne vom
Aufruf des Workflow-Management-Systems bis zu dessen Beendigung
bezeichnet. Sie umfasst damit alle Vorgänge eines Laufs, die
Ausführung der Rechenschritte ebenso wie die Vorbereitung und die
Übergaben zwischen den Schritten, und entspricht der Zeit, die
Nutzende bei der Ausführung eines Workflows beobachten.

\emph{Overhead:} Der Overhead bezeichnet die zusätzliche Zeit, die durch den Einsatz eines Workflow-Management-Systems gegenüber der Ausführung der reinen Rechenlogik entsteht. Er wird als Differenz zwischen der Gesamtlaufzeit eines Laufs und der Laufzeit eines Referenzlaufs bestimmt, der dieselben Rechenschritte in derselben Reihenfolge ohne Workflow-Management-System ausführt. Berichtet wird der Median der je Wiederholung gebildeten
Differenzen zwischen System- und Referenzlauf.

\emph{Übergangslatenz:} Die Übergangslatenz bezeichnet die
Zeitspanne zwischen dem Ende eines Rechenschritts und dem Beginn des
nachfolgenden. Sie erfasst damit den Zeitbedarf, der beim Übergang
zwischen zwei Schritten entsteht, und wird für jeden Übergang eines
Musters einzeln bestimmt.

\emph{Startspreizung:} Die Startspreizung bezeichnet den Abstand
zwischen dem Beginn der ersten und dem Beginn der letzten parallelen
Verarbeitungsinstanz eines Scatter-Gather-Laufs. Sie erfasst, wie
eng beieinander ein System die Instanzen der parallelen Phase startet,
und ist bei einem einzelnen Teilstück nicht definiert.

Die vier Größen erfassen den zeitlichen Einfluss eines
Workflow-Management-Systems auf zwei Ebenen. Gesamtlaufzeit und
Overhead beschreiben den Zeitbedarf eines vollständigen Laufs,
Übergangslatenz und Startspreizung betrachten zwei
Koordinationsvorgänge innerhalb eines Laufs gesondert. Makespan und Overhead knüpfen an bestehende Arbeiten zur
Analyse des Laufzeitverhaltens wissenschaftlicher Workflows an.
Chen und Deelman analysieren den Zeitbedarf eines Workflow-Laufs
anhand des Makespan und verschiedener Overhead-Anteile des
Workflow-Management-Systems und der Ausführungsumgebung~\cite{chen2011overhead}. Die beiden
Koordinationsgrößen werden für diese Arbeit eingeführt. Sie bestimmen den zeitlichen Abstand zwischen zwei aufeinanderfolgenden
Rechenschritten beziehungsweise die Spreizung der Startzeitpunkte
paralleler Instanzen. Beide Größen lassen sich in allen fünf
Systemen auf dieselbe Weise erheben, ohne auf systemeigene
Protokolle zurückzugreifen.

\section{Benchmark-Workflows und Versuchsplan}
\label{sec:methodik:versuchsplan}

Für die Untersuchung wurden zwei Benchmark-Workflows entworfen, die
die beiden in Kapitel~\ref{sec:basics} eingeführten Workflow-Muster
umsetzen. Der Pipeline-Workflow bildet eine vollständig
sequenzielle Verarbeitungskette ab, der Scatter-Gather-Workflow
erweitert diese Kette um eine parallele Verarbeitungsphase mit
anschließender Zusammenführung. Die Wahl fiel auf diese beiden
Muster, da sie in der Literatur als grundlegende Strukturen
wissenschaftlicher Workflows beschrieben werden. Suter et al. nennen Datenverarbeitungs-Pipelines und Fan-out/Fan-in-Ausführungsmuster als einfache Formen wissenschaftlicher Workflows~\cite[S.~4]{suter2026terminology}.

Der Pipeline-Workflow besteht aus vier Schritten.
\texttt{generate\_input} erzeugt einen reproduzierbaren
Eingabedatensatz. \texttt{preprocess} bereitet die Daten für die
Verarbeitung vor, indem Werte unterhalb einer Schwelle entfernt
werden. \texttt{compute} führt als einziger Schritt die
rechenintensive Verarbeitung aus. \texttt{postprocess} fasst das
Ergebnis zu einer abschließenden Ausgabedatei zusammen.

Der Scatter-Gather-Workflow übernimmt die Vor- und
Nachbereitungsschritte des Pipeline-Workflows und erweitert die
Verarbeitung um eine parallele Phase. Nach der Vorverarbeitung
zerlegt \texttt{split} den Datensatz in eine einstellbare Zahl von
Teilstücken. Für jedes Teilstück wird eine eigene Instanz von
\texttt{compute} ausgeführt. \texttt{aggregate} führt die
Teilergebnisse zu einem gemeinsamen Ergebnis zusammen, bevor
\texttt{postprocess} den Workflow abschließt.

Abbildung~\ref{fig:methodik:workflows} zeigt den Aufbau beider
Workflows.

\begin{figure}[H]
  \centering
  \begin{tikzpicture}[
      node distance=6.5mm and 8mm,
      wfstep/.style={draw=ctcolorgraylight, fill=ctcolorgraylighter!45,
        rounded corners=2pt, minimum width=26mm, minimum height=6mm,
        align=center, font={\small\ttfamily}},
      wfload/.style={wfstep, draw=ctcolormain!70, fill=ctcolormain!8},
      wfinst/.style={wfload, minimum width=18mm},
      lbl/.style={font=\scriptsize\itshape, text=ctcolorgray, anchor=west,
        xshift=1.5mm},
      head/.style={font=\small\bfseries, text=ctcolorgray},
      arr/.style={-stealth, semithick, draw=ctcolorgray,
        shorten <=0.5pt, shorten >=0.5pt}]
    % Pipeline (links)
    \node[head] (h1) {Pipeline};
    \node[wfstep, below=4mm of h1] (g1) {generate\_input};
    \node[wfstep, below=of g1] (p1) {preprocess};
    \node[wfload, below=of p1] (c1) {compute};
    \node[wfstep, below=of c1] (o1) {postprocess};
    \node[lbl] at (g1.east) {Datenerzeugung};
    \node[lbl] at (p1.east) {Filterung};
    \node[lbl] at (c1.east) {Matrixberechnung};
    \node[lbl] at (o1.east) {Zusammenfassung};
    \draw[arr] (g1) -- (p1);
    \draw[arr] (p1) -- (c1);
    \draw[arr] (c1) -- (o1);
    % Scatter-Gather (rechts)
    \node[head, right=58mm of h1] (h2) {Scatter-Gather};
    \node[wfstep, below=4mm of h2] (g2) {generate\_input};
    \node[wfstep, below=of g2] (p2) {preprocess};
    \node[wfstep, below=of p2] (s2) {split};
    \node[wfinst, below left=8.5mm and 2mm of s2] (ca) {compute$_1$};
    \node[wfinst, below right=8.5mm and 2mm of s2] (cb) {compute$_n$};
    \node[text=ctcolorgray] at ($(ca.east)!0.5!(cb.west)$) {$\dots$};
    \node[wfstep, below=23mm of s2] (a2) {aggregate};
    \node[wfstep, below=of a2] (o2) {postprocess};
    \node[lbl] at (s2.east) {Zerlegung in $n$ Teilstücke};
    \node[lbl] at (a2.east) {Zusammenführung};
    \draw[arr] (g2) -- (p2);
    \draw[arr] (p2) -- (s2);
    \draw[arr] (s2.south) -- (ca.north);
    \draw[arr] (s2.south) -- (cb.north);
    \draw[arr] (ca.south) -- (a2.north west);
    \draw[arr] (cb.south) -- (a2.north east);
    \draw[arr] (a2) -- (o2);
  \end{tikzpicture}
  \caption{Aufbau der beiden Benchmark-Workflows. Hervorgehoben ist der \texttt{compute}-Schritt als einziger Träger der Rechenlast, beim Scatter-Gather-Muster in $n$ parallelen Instanzen.}
  \label{fig:methodik:workflows}
\end{figure}

Die Struktur der beiden Workflows orientiert sich an
Grundstrukturen wissenschaftlicher Workflows, die Bharathi et al.
beschreiben, darunter sequenzielle Verarbeitungsketten, die
Zerlegung von Datensätzen in Teilstücke und deren anschließende
Zusammenführung~\cite[S.~3]{bharathi2008characterization}.
Die konkrete Ausgestaltung der einzelnen Schritte stellt eine
vereinfachte Benchmark-Implementierung dieser Grundstrukturen dar.

Alle Schritte sind Python-Skripte, die ihre Ein- und Ausgabepfade
als Argumente entgegennehmen und ihre Ergebnisse als Dateien
ablegen. \texttt{generate\_input} erzeugt einhundert ganzzahlige
Werte mit einem festen Startwert des Zufallszahlengenerators, die
Eingabe ist dadurch in jedem Lauf identisch. Die eigentliche
Rechenlast liegt ausschließlich im \texttt{compute}-Schritt: Für
jeden Eingabewert wird das Produkt zweier Matrizen der Größe
$96 \times 96$ berechnet, die deterministisch aus dem Wert
abgeleitet werden, und zu einer Prüfsumme zusammengefasst. Die
übrigen Schritte beschränken sich auf das Erzeugen, Filtern,
Zerlegen und Zusammenführen der Textdateien und tragen keine
nennenswerte Rechenlast.

Für die Erzeugung von Workflow-Benchmarks existieren auch Werkzeuge
wie WfBench, das realitätsnahe Rechenlasten und Workflow-Strukturen
erzeugt~\cite[S.~101]{wfbench2022}. Für diese Arbeit wurden
stattdessen bewusst einfache Benchmark-Workflows verwendet, um die
Rechenlogik und die Versuchsbedingungen über alle fünf Systeme
hinweg möglichst kontrolliert und einheitlich zu halten.

Die Untersuchung verwendet drei Rechenlasten, die sich
ausschließlich in der Zahl der Wiederholungen der Matrixberechnung
je Eingabewert unterscheiden: einmal bei der kurzen, dreizehnmal
bei der mittleren und sechzigmal bei der langen Rechenlast. Die
Ausführung der reinen Rechenlogik benötigt auf der lokalen Ebene
etwa fünf Sekunden, eine Minute beziehungsweise viereinhalb
Minuten. Die drei Rechenlasten ermöglichen damit die Untersuchung
sowohl kurzer Workflows, bei denen der Koordinationsaufwand der
Systeme einen vergleichsweise hohen Anteil an der Gesamtlaufzeit
hat, als auch längerer Workflows, bei denen die Rechenzeit
dominiert. Beim Scatter-Gather-Workflow wird die Zahl der
Teilstücke zwischen eins, zwei und vier variiert, die parallele
Phase umfasst damit bis zu vier gleichzeitig laufende
Verarbeitungsinstanzen.

Aus den Mustern, Rechenlasten und Teilstückzahlen ergibt sich der
Versuchsplan: zwölf Konfigurationen je System, drei für das Pipeline-Muster und neun für das Scatter-Gather-Muster. Jede Konfiguration wurde bei der kurzen und der mittleren Rechenlast fünfmal wiederholt, bei der langen Rechenlast aufgrund des deutlich höheren Zeitbedarfs je Lauf dreimal. Damit stellt die Zahl der Wiederholungen einen Kompromiss zwischen Messaufwand und praktischer Durchführbarkeit dar.


\section{Messverfahren}
\label{sec:methodik:messverfahren}

Die Messung erfolgt auf zwei Ebenen. Eine äußere Messung um den
Lauf liefert die Gesamtlaufzeit und, im Vergleich mit einem
Referenzlauf, den Overhead. Eine innere Instrumentierung in den
Benchmark-Skripten liefert die Zeitpunkte, aus denen Übergangslatenz
und Startspreizung berechnet werden.
Abbildung~\ref{fig:methodik:messung} gibt einen Überblick über das
Mess- und Auswertungsverfahren einer Konfiguration, die folgenden
Abschnitte beschreiben die einzelnen Bestandteile.

\begin{figure}[H]
  \centering
  \begin{tikzpicture}[
      node distance=4mm,
      box/.style={draw=ctcolorgraylight, fill=ctcolorgraylighter!45,
        rounded corners=2pt, minimum width=62mm, minimum height=6.5mm,
        align=center, font=\small},
      clk/.style={box, fill=white, minimum height=5.5mm, font=\footnotesize},
      wfout/.style={box, minimum width=24mm, minimum height=5.5mm, font=\footnotesize},
      res/.style={box, draw=ctcolormain!70, fill=ctcolormain!8,
        minimum width=72mm, font=\small},
      grp/.style={draw=ctcolorgray, dashed, rounded corners=3pt,
        inner sep=2.5mm},
      sgrp/.style={draw=ctcolorgraylight, rounded corners=3pt,
        inner sep=2mm, inner ysep=5.5mm},
      ttl/.style={font=\small\bfseries, text=ctcolorgray},
      arr/.style={-stealth, semithick, draw=ctcolorgray}]
    \node[box] (cfg) {Konfiguration: Muster, Rechenlast, Teilstückzahl $n$};
    \node[clk, below=21mm of cfg] (r1) {äußere Messung: Startzeitpunkt};
    \node[box, below=of r1] (r2) {Schritte des Musters ohne\\ Workflow-Management-System};
    \node[clk, below=of r2] (r3) {äußere Messung: Endzeitpunkt};
    \node[sgrp, fit=(r1)(r2)(r3)] (gr) {};
    \node[wfout, right=7mm of gr] (tref) {$T_{\mathrm{Ref},r}$};
    \node[clk, below=11mm of gr.south] (s1) {äußere Messung: Startzeitpunkt};
    \node[box, below=of s1] (s2) {Schritte des Musters mit\\ Workflow-Management-System};
    \node[clk, below=of s2] (s3) {äußere Messung: Endzeitpunkt};
    \node[sgrp, fit=(s1)(s2)(s3)] (gs) {};
    \node[wfout, right=7mm of gs] (tsys) {$T_{\mathrm{Sys},r}$};
    \node[box, below=9mm of gs.south] (tim) {Timing-Dateien beider Läufe};
    \node[wfout, right=7mm of tim] (tco) {$L_{i,r}$, $S_r$};
    \node[grp, fit=(gr)(gs)(tim)(tref)(tco)] (gw) {};
    \node[ttl, anchor=south east, yshift=1mm] at (gw.north east) {Wiederholung $r = 1, \dots, R$};
    \node[ttl, font=\footnotesize\bfseries, anchor=north west, xshift=2mm, yshift=-1.2mm] at (gr.north west) {Referenzlauf};
    \node[ttl, font=\footnotesize\bfseries, anchor=north west, xshift=2mm, yshift=-1.2mm] at (gs.north west) {Systemlauf};
    \node[box, below=10mm of gw] (agg) {Median über die Wiederholungen};
    \node[res, below=7mm of agg] (o1)
  {Gesamtlaufzeit: Median von $T_{\mathrm{Sys},r}$};
\node[res, below=4mm of o1] (o2)
  {Overhead: Median von $T_{\mathrm{Sys},r}-T_{\mathrm{Ref},r}$};
\node[res, below=4mm of o2] (o3)
  {Übergangslatenz: Median von $L_{i,r}$, Startspreizung: Median von $S_r$};
    \draw[arr] (cfg) -- (r1);
    \draw[arr] (r1) -- (r2);
    \draw[arr] (r2) -- (r3);
    \draw[arr] (gr.east) -- (tref);
    \draw[arr] (gr.south) -- (gs.north);
    \draw[arr] (s1) -- (s2);
    \draw[arr] (s2) -- (s3);
    \draw[arr] (gs.east) -- (tsys);
    \draw[arr] (gs.south) -- (tim);
    \draw[arr] (tim.east) -- (tco);
    \draw[arr] (gw.south) -- (agg);
    \draw[arr] (agg) -- (o1);
    \draw[arr] (o1) -- (o2);
    \draw[arr] (o2) -- (o3);
  \end{tikzpicture}
  \caption{Ablauf des Mess- und Auswertungsverfahrens für eine Konfiguration. Jede Wiederholung umfasst einen Referenzlauf und den zugehörigen Systemlauf, beide werden nach demselben Verfahren von außen gemessen und erzeugen Timing-Dateien. Aus den Wiederholungen werden anschließend die berichteten Kennwerte gebildet.}
  \label{fig:methodik:messung}
\end{figure}

\subsection{Referenzmessung}
\label{sec:methodik:referenz}

Der Overhead eines Systems lässt sich nicht unmittelbar messen,
sondern nur als Differenz zu einer Ausführung derselben Rechenlogik
ohne Workflow-Management-System. Zu jeder Wiederholung einer
Konfiguration wurde deshalb neben dem Systemlauf ein Referenzlauf
erhoben, der als Bezugsgröße dient.

Der Referenzlauf verwendet dieselben Benchmark-Skripte und
dieselben fachlichen Ein- und Ausgaben wie der zugehörige
Systemlauf und führt sie in der Reihenfolge des jeweiligen Musters
aus. Referenz- und Systemläufe werden für jedes System durch ein
Steuerskript gestartet und gemessen. Für den Referenzlauf startet
dieses Steuerskript die Benchmark-Skripte unmittelbar als eigene
Prozesse. Beim Pipeline-Muster laufen die vier Schritte
nacheinander ab. Beim Scatter-Gather-Muster folgen auf die
Vorverarbeitung und die Zerlegung so viele gleichzeitig gestartete
Verarbeitungsprozesse, wie Teilstücke vorgesehen sind, anschließend
das Zusammenführen und die Nachbereitung. Da die Eingabe über einen
festen Startwert erzeugt wird, verarbeiten Referenz- und Systemlauf
dieselben Werte. Der Referenzlauf bildet damit dieselbe Rechenlogik
sowie dieselbe Abfolge und Parallelität der Rechenschritte ab wie
der Systemlauf, jedoch ohne die Steuerung durch ein
Workflow-Management-System.

Innerhalb einer Wiederholung wurde stets zuerst der Referenzlauf
und anschließend der zugehörige Systemlauf ausgeführt. Beide liefen
in eigenen Arbeitsverzeichnissen und wurden nach demselben
Verfahren von außen gemessen. Durch die unmittelbare zeitliche
Paarung bleiben Unterschiede der Ausführungsbedingungen gering, der
wesentliche Unterschied zwischen beiden Läufen liegt in der
Steuerung durch das Workflow-Management-System.

\subsection{Äußere Messung}
\label{sec:methodik:aussen}

Die äußere Messung erfasst die Gesamtlaufzeit eines Laufs. Das
Steuerskript nimmt unmittelbar vor dem Start eines Laufs und nach
dessen vollständigem Abschluss jeweils einen Zeitstempel. Die
Zeitspanne dazwischen bildet die Gesamtlaufzeit
$T_\mathrm{Sys}$ eines Systemlaufs beziehungsweise
$T_\mathrm{Ref}$ des zugehörigen Referenzlaufs.

Der Beginn der Messung liegt bei Nextflow, StreamFlow, Merlin und
AiiDA unmittelbar vor dem Startbefehl. Bei Pegasus beginnt die
Messung bereits vor der Planung, sodass auch die Planungszeit in
der Gesamtlaufzeit enthalten ist. Wie das Ende eines Laufs festgestellt wird, hängt vom
Ausführungsmodell ab. Bei Nextflow und StreamFlow kehrt der Aufruf
erst zurück, wenn der Workflow abgeschlossen ist, das Ende der
Messung fällt mit der Rückkehr des Aufrufs zusammen. Merlin, AiiDA
und Pegasus übergeben die Ausführung an eine eigenständige
Komponente, der Startbefehl kehrt vorher zurück. Das Steuerskript
fragt in diesen Fällen den Zustand des Laufs zyklisch ab und
beendet die Messung, sobald alle Schritte abgeschlossen sind.
Tabelle~\ref{tab:methodik:endeerkennung} fasst je System zusammen,
womit ein Lauf gestartet und woran sein Ende erkannt wird.

\begin{table}[H]
  \caption{Start der Läufe und Erkennung des Laufendes in der äußeren Messung}
  \label{tab:methodik:endeerkennung}
  \centering
  \footnotesize
  \renewcommand{\arraystretch}{1.3}
  \begin{tabular}{>{\raggedright\arraybackslash}p{1.8cm}
                  >{\raggedright\arraybackslash}p{3.4cm}
                  >{\raggedright\arraybackslash}p{7.0cm}}
    \hline
    \textbf{System} & \textbf{Start der Messung} & \textbf{Erkennung des Laufendes} \\
    \hline
    Nextflow & vor \texttt{nextflow run} & Rückkehr des blockierenden Aufrufs \\
    StreamFlow & vor \texttt{streamflow run} & Rückkehr des blockierenden Aufrufs \\
    Merlin & vor \texttt{merlin run} & \texttt{MERLIN\_STATUS.json} mit Status \texttt{FINISHED} und \texttt{MERLIN\_FINISHED} im Workspace des terminalen Schritts \texttt{postprocess}, Abfrage alle 20\,ms \\
    AiiDA & vor \texttt{submit()} der WorkChain & terminaler Zustand der obersten WorkChain, Abfrage alle 50\,ms \\
    Pegasus & vor \texttt{pegasus-plan} & erfolgreicher Workflow-Zustand über \texttt{pegasus-status -l}, Abfrage alle 500\,ms \\
    \hline
  \end{tabular}
\end{table}

Bei Pegasus umfasst die Messung damit den vorgelagerten
Planungsschritt, da dieser Teil der Ausführung eines Workflows ist. Dauerhaft laufende Komponenten wie Worker, Daemon oder HTCondor sind beim Start der Messung bereits aktiv, ihr Start zählt nicht zur Gesamtlaufzeit. Die äußere Messung umfasst damit in allen Systemen denselben Abschnitt: Sie beginnt mit der Übergabe des vorbereiteten Workflows an das jeweilige System und endet, sobald dieses den vollständigen Workflow als abgeschlossen meldet.

\subsection{Innere Instrumentierung}
\label{sec:methodik:instrumentierung}

Die beiden Koordinationsgrößen erfordern Zeitpunkte einzelner
Schritte innerhalb eines Laufs. Diese Zeitpunkte werden nicht den
Protokollen der Systeme entnommen, sondern in den
Benchmark-Skripten selbst erhoben. Da alle fünf Systeme dieselben
Skripte ausführen, entsteht dadurch in jedem System dieselbe
Datenquelle, unabhängig davon, welche Aufzeichnungen ein System von
sich aus führt und in welcher Form und Auflösung es sie ablegt.

Jeder Schritt nimmt zu Beginn seiner Ausführung und nach Abschluss
der eigentlichen Verarbeitung jeweils einen Zeitstempel und legt anschließend eine Timing-Datei
in seinem Arbeitsverzeichnis ab, die den Namen des Schritts sowie
beide Zeitpunkte enthält. Aus den Zeitpunkten aufeinanderfolgender
Schritte ergibt sich die Übergangslatenz, aus den Startzeitpunkten
der parallelen Instanzen eines Scatter-Gather-Laufs die
Startspreizung. Die Instanzen tragen dabei die Nummer ihres
Teilstücks im Namen und lassen sich dadurch einzeln auswerten.
Listing~\ref{lst:methodik:timing} zeigt den Aufbau eines
instrumentierten Schritts.

\begin{lstlisting}[language=Python,
  caption={Aufbau eines instrumentierten Benchmark-Schritts},
  label={lst:methodik:timing}]
task_start_ns = time.time_ns()

# fachliche Verarbeitung des Schritts

task_end_ns = time.time_ns()
write_timing("preprocess", task_start_ns, task_end_ns)


def write_timing(task_name, start_ns, end_ns):
    Path(f"timing_{task_name}.txt").write_text(
        f"task_name={task_name}\n"
        f"task_start_ns={start_ns}\n"
        f"task_end_ns={end_ns}\n"
    )
\end{lstlisting}

Als Zeitquelle dient die Systemuhr in Nanosekunden
(\texttt{time.time\_ns()}). Da die Übergangslatenzen aus
Zeitpunkten verschiedener Prozesse berechnet werden, müssen alle
Schritte dieselbe Zeitbasis verwenden. Die gewählte Zeitquelle
ermöglicht dadurch den direkten Vergleich der Zeitstempel aller
Benchmark-Schritte.

Da Referenzlauf und Systemlauf dieselben Benchmark-Skripte
verwenden, entstehen in beiden Fällen Timing-Dateien mit
identischem Aufbau. Die Instrumentierung ist dadurch unabhängig
vom verwendeten Workflow-Management-System und liefert für alle
fünf Systeme dieselbe Grundlage zur Berechnung der
Koordinationsgrößen.

\subsection{Auswertung}
\label{sec:methodik:auswertung}

Nach Abschluss aller Wiederholungen einer Konfiguration werden die
äußeren Laufzeiten und die Timing-Dateien zusammengeführt. Aus den
äußeren Laufzeiten ergeben sich Gesamtlaufzeit und Overhead, aus
den Timing-Dateien die beiden Koordinationsgrößen. Anschließend
werden die Messwerte über die Wiederholungen zusammengefasst.

Sei $R$ die Zahl der Wiederholungen einer Konfiguration und
$r \in \{1,\ldots,R\}$ der Index einer Wiederholung,
$T_{\mathrm{Sys},r}$ und $T_{\mathrm{Ref},r}$ die in Wiederholung
$r$ von außen gemessenen Laufzeiten des Systemlaufs und des
zugehörigen Referenzlaufs. Die berichtete Gesamtlaufzeit $T$ und
der Overhead $O$ sind dann
\[
  T = \operatorname{Median}_{r}\bigl(T_{\mathrm{Sys},r}\bigr),
  \qquad
  O = \operatorname{Median}_{r}\bigl(T_{\mathrm{Sys},r}
      - T_{\mathrm{Ref},r}\bigr).
\]
Der Overhead wird damit entsprechend seiner Definition in
Abschnitt~\ref{sec:methodik:messgroessen} als Median der je
Wiederholung gebildeten Differenzen bestimmt.

Die Koordinationsgrößen ergeben sich aus den Timing-Dateien. Für
einen Übergang $i$ vom Schritt $a$ zum Schritt $b$ bezeichnet
$L_{i,r}$ den Abstand zwischen dem Ende von $a$ und dem Beginn von
$b$, für die parallele Phase eines Scatter-Gather-Laufs mit $n$
Teilstücken bezeichnet $S_r$ den Abstand zwischen dem frühesten und
dem spätesten Start der Instanzen:
\[
  L_{i,r} = t^{\text{start}}_{b,r} - t^{\text{ende}}_{a,r},
  \qquad
  S_r = \max_{1 \le k \le n} t^{\text{start}}_{k,r}
      - \min_{1 \le k \le n} t^{\text{start}}_{k,r}.
\]
Berichtet werden auch hier die Mediane über die Wiederholungen,
$L_i = \operatorname{Median}_{r}(L_{i,r})$ und
$S = \operatorname{Median}_{r}(S_r)$.

Der Median wurde dem arithmetischen Mittel vorgezogen, da bei einer
kleinen Zahl von Wiederholungen vereinzelte ungewöhnlich lange
Laufzeiten das Ergebnis sonst bestimmen würden.

Die Steuerskripte legen die Messwerte der einzelnen Wiederholungen
als Rohdaten ab und erzeugen aus den Timing-Dateien die
Koordinationswerte. Für die endgültige Auswertung werden die
äußeren Laufzeitmessungen von einem separaten Auswertungsskript
eingelesen und nach dem oben beschriebenen Verfahren
zusammengefasst. Die bereits berechneten Koordinationswerte werden
unverändert übernommen.

Auf der HPC-Ebene bleiben Instrumentierung und Auswertung
unverändert. Die Ausführungsbedingungen unterscheiden sich dort
jedoch, weshalb die Ergebnisse beider Ebenen getrennt betrachtet
werden. Kapitel~\ref{sec:mogon} geht darauf ein.

\subsection{Grenzen des Messverfahrens}
\label{sec:methodik:grenzen}

Das Messverfahren ist darauf ausgerichtet, die
Workflow-Management-Systeme unter vergleichbaren Bedingungen und
mit derselben Rechenlogik zu vermessen. Einige Eigenschaften des
Versuchsaufbaus sind bei der Einordnung der Ergebnisse zu
berücksichtigen.

Die Zahl der Wiederholungen beträgt für die kurze und die mittlere
Rechenlast fünf, für die lange Rechenlast aufgrund des deutlich
höheren Zeitbedarfs drei. Die berichteten Mediane beruhen damit je
nach Rechenlast auf unterschiedlich vielen Beobachtungen.

Referenzlauf und Systemlauf werden innerhalb einer Wiederholung
unmittelbar nacheinander ausgeführt und bilden dadurch möglichst
vergleichbare Bedingungen. Dennoch unterliegen beide Läufe den
üblichen Schwankungen der Ausführungsumgebung, die sich nicht
vollständig vermeiden lassen.

Das Messverfahren erfasst die von außen beobachtbaren Laufzeiten
eines Systems. Die Messgrößen beschreiben damit den zeitlichen
Effekt der Ausführung, erlauben jedoch keine unmittelbare Zuordnung
der gemessenen Zeiten zu einzelnen internen Mechanismen oder
Komponenten eines Workflow-Management-Systems. Die Auswirkungen
dieser Einschränkungen auf die Interpretation der Ergebnisse werden
in Abschnitt~\ref{sec:diskussion:grenzen} diskutiert.